\documentclass[11pt]{article}
% ======================================================================
%  examstyle.tex — EPFL-style exam layout for the weekly Time Series exams
%  Shared by every Exam_Week_NN(.tex / _Solutions.tex). Compiles with pdflatex.
% ======================================================================
\usepackage[utf8]{inputenc}
\usepackage[T1]{fontenc}
\usepackage[english]{babel}
\usepackage[a4paper,margin=2.4cm]{geometry}
\usepackage{amsmath,amssymb,amsthm,mathtools}
\usepackage{bm}
\usepackage{enumitem}
\usepackage{array}

\setlength{\parindent}{0pt}
\setlength{\parskip}{0.5em}
\emergencystretch=3em
\allowdisplaybreaks[1]

% ---------- math macros (same names as the rest of the project) ----------
\newcommand{\E}{\mathbb{E}}
\DeclareMathOperator{\Var}{Var}
\DeclareMathOperator{\Cov}{Cov}
\DeclareMathOperator{\Corr}{Corr}
\DeclareMathOperator*{\argmin}{arg\,min}
\DeclareMathOperator{\MSE}{MSE}
\DeclareMathOperator{\trace}{tr}
\newcommand{\Reals}{\mathbb{R}}
\newcommand{\Ints}{\mathbb{Z}}
\newcommand{\Comp}{\mathbb{C}}
\newcommand{\Nat}{\mathbb{N}}
\newcommand{\Prob}{\mathbb{P}}
\newcommand{\Tr}{^{\mathsf{T}}}
\newcommand{\conj}[1]{\overline{#1}}
\newcommand{\abs}[1]{\left\lvert #1 \right\rvert}
\newcommand{\norm}[1]{\left\lVert #1 \right\rVert}
\newcommand{\ind}{\mathbf{1}}
\newcommand{\eps}{\varepsilon}
\newcommand{\diff}{\nabla}
\newcommand{\acvs}{\gamma}
\newcommand{\acf}{\rho}
\newcommand{\pacf}{\alpha}
\newcommand{\sdf}{\mathcal{S}}
\newcommand{\Bsh}{B}
\newcommand{\iid}{\overset{\text{iid}}{\sim}}

\setlist[enumerate]{leftmargin=2em,topsep=2pt,itemsep=4pt}
\setlist[itemize]{leftmargin=1.6em,topsep=2pt,itemsep=2pt}

\newcounter{exo}

% ---------- exam cover header :  \examheader{exam title}{duration} ----------
\newcommand{\examheader}[2]{%
  \thispagestyle{empty}
  \begin{center}
    {\Large\bfseries Time Series}\\[3pt]
    Spring Semester 2026, EPFL\\
    \'Ecole Polytechnique F\'ed\'erale de Lausanne\\
    Prof.\ Dr.\ Sofia C.\ Olhede\\[7pt]
    \hrule height 0.8pt
    \vspace{9pt}
    {\large\bfseries Practice Exam --- #1}\\[4pt]
    {\normalsize Duration: #2 \quad$\cdot$\quad No calculator \quad$\cdot$\quad Answer on paper}\\
    \vspace{8pt}
    \hrule height 0.8pt
  \end{center}
  \vspace{14pt}
  \begin{tabular}{@{}p{8.2cm}@{}p{8cm}@{}}
    Name:\enspace\hrulefill & Surname:\enspace\hrulefill
  \end{tabular}\\[14pt]
  SCIPER (Student Card Number):\enspace\hrulefill
  \vspace{16pt}
}

% ---------- points table (fixed: 4 problems) ----------
\newcommand{\pointstable}{%
  \begin{center}
  \renewcommand{\arraystretch}{1.5}
  \begin{tabular}{|p{5cm}|p{4cm}|}
    \hline
    \textbf{Exercise} & \textbf{Points}\\\hline
    1 & \\\hline
    2 & \\\hline
    3 & \\\hline
    4 & \\\hline
    \textbf{Total} & \\\hline
  \end{tabular}
  \end{center}
}

% ---------- problem heading (exam: one problem per page) ----------
\newcommand{\problem}[1]{%
  \clearpage
  \stepcounter{exo}%
  \begin{center}\textbf{\large Problem \theexo\ (#1 points)}\end{center}
  \vspace{4pt}
}
% ---------- answer space: fill the statement page + one blank answer page ----------
% Put \answerpage at the END of each problem so every exercise gets its own page(s).
\newcommand{\answerpage}{\par\vfill\newpage\null\newpage}

% ---------- solutions document ----------
\newcommand{\solheader}[1]{%
  \begin{center}
    {\Large\bfseries Time Series --- Solutions}\\[3pt]
    {\large Practice Exam --- #1}\\[2pt]
    EPFL, MATH-342 \quad$\cdot$\quad Prof.\ S.\ C.\ Olhede\\[5pt]
    \hrule height 0.8pt
  \end{center}
  \vspace{10pt}
}
\newcommand{\solproblem}[1]{%
  \bigskip
  \stepcounter{exo}%
  {\large\bfseries Problem \theexo\ (#1 points)}\par\nobreak\vspace{2pt}
}
% Rappel de l'énoncé avant la solution (dans les corrigés)
\newcommand{\statementstart}{\par\vspace{1pt}{\bfseries\itshape Statement.}\par\nobreak\begingroup\small\itshape}
\newcommand{\statementend}{\par\endgroup\vspace{5pt}\hrule\vspace{6pt}{\bfseries Solution.}\par\nobreak\vspace{2pt}}

% --- local macros not in examstyle ---
\DeclareMathOperator{\diag}{diag}
\begin{document}

\solheader{Week 10: Forecasting}

% ---------------------------------------------------------------
\solproblem{5}
\statementstart
\textbf{(Forecasting a causal AR(1): point forecast, error variance, interval.)}
Let $\{X_t\}$ be the zero-mean Gaussian AR(1) process
\[
X_t=\phi X_{t-1}+\eps_t,\qquad \abs\phi<1,\quad \eps_t\iid\mathcal N(0,\sigma^2),
\]
and write $X^{\,n}_{n+h}$ for the best linear predictor of $X_{n+h}$ based on
$X_1,\dots,X_n$, with $h$-step forecast error $e_n(h)=X_{n+h}-X^{\,n}_{n+h}$.

\textbf{(a)} Give the $\mathrm{MA}(\infty)$ weights $\{\psi_j\}$ of the process. Using that the process is Gaussian (so the best linear predictor equals the conditional expectation), show that the point forecast is $X^{\,n}_{n+h}=\phi^{\,h}X_n$, and state its limit as $h\to\infty$. \hfill(1.5 pts)

\textbf{(b)} Show that the $h$-step forecast-error variance is
\[
\Var\bigl(e_n(h)\bigr)=\sigma^2\,\frac{1-\phi^{2h}}{1-\phi^2},
\]
and verify that $\Var(e_n(h))\to\acvs_0$ (the marginal variance) as $h\to\infty$. \hfill(2 pts)

\textbf{(c)} Now fix $\phi=0.6$ and $\sigma^2=4$. Report $X^{\,n}_{n+2}$ and the numerical forecast-error variances $\Var(e_n(1)),\Var(e_n(2))$, then give the $95\%$ Gaussian prediction interval for $X_{n+2}$ (use $z_{0.975}=1.96$). \hfill(1.5 pts)
\statementend

\textbf{(a)} For the causal AR(1) the $\mathrm{MA}(\infty)$ representation is obtained by inverting $1-\phi B$:
\[
X_t=\frac{1}{1-\phi B}\,\eps_t=\sum_{j=0}^\infty\phi^{\,j}\eps_{t-j},
\qquad\text{so}\qquad \psi_j=\phi^{\,j}\ \ (j\ge0).
\]
Because the process is Gaussian, the best linear predictor coincides with the conditional expectation $\E[X_{n+h}\mid X_1,\dots,X_n]$. Iterating the recursion $X_{n+h}=\phi X_{n+h-1}+\eps_{n+h}$ and taking conditional expectations (every future innovation $\eps_{n+1},\dots,\eps_{n+h}$ has conditional mean $0$, being uncorrelated with the data),
\[
\begin{aligned}
X^{\,n}_{n+h}
&=\E[\phi X_{n+h-1}+\eps_{n+h}\mid X_1,\dots,X_n]
=\phi\,\E[X_{n+h-1}\mid X_1,\dots,X_n]\\
&=\phi^2\,\E[X_{n+h-2}\mid X_1,\dots,X_n]=\cdots=\phi^{\,h}X_n .
\end{aligned}
\]
Since $\abs\phi<1$, $X^{\,n}_{n+h}=\phi^{\,h}X_n\xrightarrow[h\to\infty]{}0$, the process mean: the forecast \emph{reverts to the mean}.

\textbf{(b)} By the $\psi$-weight theorem the forecast error is the sum of the unrealised future innovations,
\[
e_n(h)=X_{n+h}-X^{\,n}_{n+h}=\sum_{j=0}^{h-1}\psi_j\eps_{n+h-j}
=\sum_{j=0}^{h-1}\phi^{\,j}\eps_{n+h-j}.
\]
The innovations are uncorrelated with common variance $\sigma^2$, so
\[
\Var\bigl(e_n(h)\bigr)=\sigma^2\sum_{j=0}^{h-1}\psi_j^2
=\sigma^2\sum_{j=0}^{h-1}\phi^{2j}
=\sigma^2\,\frac{1-(\phi^2)^h}{1-\phi^2}
=\sigma^2\,\frac{1-\phi^{2h}}{1-\phi^2},
\]
using the finite geometric sum $\sum_{j=0}^{h-1}r^j=(1-r^h)/(1-r)$ with $r=\phi^2$. As $h\to\infty$, $\phi^{2h}\to0$ (since $\abs\phi<1$), so
\[
\Var\bigl(e_n(h)\bigr)\longrightarrow\frac{\sigma^2}{1-\phi^2}=\acvs_0,
\]
the marginal variance of the AR(1). The forecast-error variance increases monotonically in $h$ and saturates at $\acvs_0$: forecasting far ahead is no better than quoting the unconditional distribution.

\textbf{(c)} With $\phi=0.6$ and $\sigma^2=4$:
\[
X^{\,n}_{n+2}=\phi^2 X_n=0.36\,X_n .
\]
The forecast-error variances are
\[
\Var(e_n(1))=\sigma^2\psi_0^2=\sigma^2=4,
\qquad
\Var(e_n(2))=\sigma^2(1+\phi^2)=4(1+0.36)=4\cdot1.36=5.44 .
\]
(For reference $\acvs_0=4/(1-0.36)=4/0.64=6.25$, and indeed $4<5.44<6.25$.) A $95\%$ Gaussian prediction interval for $X_{n+2}$ is
\[
X^{\,n}_{n+2}\pm z_{0.975}\sqrt{\Var(e_n(2))}
=0.36\,X_n\pm1.96\sqrt{5.44}
=0.36\,X_n\pm1.96(2.3324)
=0.36\,X_n\pm4.571 .
\]

% ---------------------------------------------------------------
\solproblem{5}
\statementstart
\textbf{(The prediction equations, and one-step AR(2) forecasts.)} Let $\{X_t\}$ be a zero-mean second-order stationary process with autocovariance $\{\acvs_\tau\}$, and let
\[
X^{\,n}_{n+h}=\sum_{j=1}^{n}\beta_j X_j
\]
be its best linear predictor of $X_{n+h}$ based on $X=(X_1,\dots,X_n)\Tr$.

\textbf{(a)} Starting from the prediction mean square error $S(\beta)=\E\bigl[(X_{n+h}-\beta\Tr X)^2\bigr]$, derive the \emph{prediction equations} $\Gamma_n\beta=\acvs_{[h]}$, where $\Gamma_n=[\acvs_{i-j}]_{i,j=1}^n$ and $\acvs_{[h]}=(\acvs_{n+h-1},\dots,\acvs_h)\Tr$, and show that this is equivalent to the orthogonality condition $\E[(X_{n+h}-X^{\,n}_{n+h})X_k]=0$ for $k=1,\dots,n$. \hfill(2 pts)

\textbf{(b)} Assuming $\Gamma_n$ invertible, deduce the closed forms
\[
X^{\,n}_{n+h}=\acvs_{[h]}\Tr\,\Gamma_n^{-1}X,
\qquad
P^{\,n}_{n+h}=\acvs_0-\acvs_{[h]}\Tr\,\Gamma_n^{-1}\acvs_{[h]} . \hfill(1.5\text{ pts})
\]

\textbf{(c)} Let $\{X_t\}$ be the causal AR(2) process $X_t=\phi_1 X_{t-1}+\phi_2 X_{t-2}+\eps_t$. Use the result of part (a) to show that the best one-step predictors based on $X_1$ and on $X_1,X_2$ are
\[
X^{\,1}_2=\frac{\acvs_1}{\acvs_0}X_1=\acf_1 X_1,
\qquad
X^{\,2}_3=\phi_1 X_2+\phi_2 X_1 .
\]
(For the second, you need \emph{not} invert $\Gamma_2$.) \hfill(1.5 pts)
\statementend

\textbf{(a)} Expand the prediction MSE as a quadratic in $\beta$, using $\Var(X)=\Gamma_n$ and $\Cov(X_{n+h},X)=\acvs_{[h]}$ (the $k$-th entry of $\Cov(X_{n+h},X)$ is $\Cov(X_{n+h},X_k)=\acvs_{n+h-k}$, which is exactly the $k$-th entry of $\acvs_{[h]}$):
\[
\begin{aligned}
S(\beta)
&=\E\bigl[(X_{n+h}-\beta\Tr X)^2\bigr]
=\Var(X_{n+h}-\beta\Tr X)\\
&=\acvs_0-2\beta\Tr\Cov(X_{n+h},X)+\beta\Tr\Var(X)\beta
=\acvs_0-2\beta\Tr\acvs_{[h]}+\beta\Tr\Gamma_n\beta .
\end{aligned}
\]
This is a quadratic in $\beta$ bounded below by $0$ ($\Gamma_n$ is positive semi-definite), so its minimiser solves $\nabla_\beta S=0$:
\[
\nabla_\beta S=-2\acvs_{[h]}+2\Gamma_n\beta=0
\qquad\Longrightarrow\qquad
\boxed{\;\Gamma_n\beta=\acvs_{[h]}.\;}
\]
\emph{Equivalence to orthogonality.} Row $k$ of $\Gamma_n\beta=\acvs_{[h]}$ reads
\[
\sum_{j=1}^n\beta_j\,\acvs_{j-k}=\acvs_{n+h-k},\qquad k=1,\dots,n .
\]
On the other hand $\E[(X_{n+h}-\sum_j\beta_jX_j)X_k]=\Cov(X_{n+h},X_k)-\sum_j\beta_j\Cov(X_j,X_k)=\acvs_{n+h-k}-\sum_j\beta_j\acvs_{j-k}$. Setting this to $0$ for each $k$ gives exactly the same $n$ equations. Hence ``error orthogonal to every observation'' $\Longleftrightarrow$ $\Gamma_n\beta=\acvs_{[h]}$: the best linear predictor is the orthogonal projection of $X_{n+h}$ onto $\mathrm{span}\{X_1,\dots,X_n\}$.

\textbf{(b)} If $\Gamma_n$ is invertible then $\beta=\Gamma_n^{-1}\acvs_{[h]}$, so
\[
X^{\,n}_{n+h}=\beta\Tr X=\bigl(\Gamma_n^{-1}\acvs_{[h]}\bigr)\Tr X=\acvs_{[h]}\Tr\Gamma_n^{-1}X
\]
(using $(\Gamma_n^{-1})\Tr=\Gamma_n^{-1}$ since $\Gamma_n$ is symmetric). For the MSE, substitute $\beta\Tr\Gamma_n\beta=\beta\Tr\acvs_{[h]}$ into $S(\beta)$:
\[
P^{\,n}_{n+h}=\acvs_0-2\beta\Tr\acvs_{[h]}+\beta\Tr\Gamma_n\beta
=\acvs_0-2\beta\Tr\acvs_{[h]}+\beta\Tr\acvs_{[h]}
=\acvs_0-\beta\Tr\acvs_{[h]}
=\acvs_0-\acvs_{[h]}\Tr\Gamma_n^{-1}\acvs_{[h]} .
\]

\textbf{(c)} \emph{Based on $X_1$ (so $n=1$, $h=1$).} The prediction equation $\Gamma_1\beta=\acvs_{[1]}$ is the scalar equation $\acvs_0\beta_1=\acvs_1$, hence $\beta_1=\acvs_1/\acvs_0=\acf_1$ and
\[
X^{\,1}_2=\beta_1X_1=\frac{\acvs_1}{\acvs_0}X_1=\acf_1 X_1 .
\]
\emph{Based on $X_1,X_2$ (so $n=2$, $h=1$).} Rather than invert $\Gamma_2$, use the model: $X_3=\phi_1X_2+\phi_2X_1+\eps_3$, so the candidate predictor $\phi_1X_2+\phi_2X_1$ has error
\[
X_3-(\phi_1X_2+\phi_2X_1)=\eps_3 .
\]
The future white noise $\eps_3$ is uncorrelated with the past observations $X_1,X_2$ (which depend only on $\eps_s$, $s\le2$), so
\[
\E\bigl[(X_3-\phi_1X_2-\phi_2X_1)X_1\bigr]=\E[\eps_3X_1]=0,
\qquad
\E\bigl[(X_3-\phi_1X_2-\phi_2X_1)X_2\bigr]=\E[\eps_3X_2]=0 .
\]
The error is orthogonal to both observations, so by the prediction-equations theorem of part (a) this candidate \emph{is} the best linear predictor:
\[
X^{\,2}_3=\phi_1X_2+\phi_2X_1 .
\]
(General fact: for an AR$(p)$ with $n\ge p$, the one-step forecast is simply $\sum_{i=1}^p\phi_iX_{n+1-i}$.)

% ---------------------------------------------------------------
\solproblem{5}
\statementstart
\textbf{(The causal-ARMA predictor via $\psi$-weights, and a worked ARMA(1,1).)}

\textbf{(a)} For a causal ARMA process with linear representation $X_t=\sum_{j=0}^\infty\psi_j\eps_{t-j}$, prove that the best linear predictor of $X_{n+h}$ based on $X_1,\dots,X_n$ is
\[
X^{\,n}_{n+h}=\sum_{j=h}^{\infty}\psi_j\,\eps_{n+h-j},
\qquad\text{with}\qquad
P^{\,n}_{n+h}=\sigma^2\sum_{j=0}^{h-1}\psi_j^2 .
\]
(\emph{Hint:} any linear predictor based on $X_1,\dots,X_n$ can be written $\sum_{j\ge0}c_j\eps_{n-j}$; expand the MSE and minimise over the $c_j$.) \hfill(2 pts)

\textbf{(b)} Consider the causal, invertible ARMA(1,1) process
\[
X_t=\tfrac12 X_{t-1}+\eps_t+\tfrac{3}{10}\eps_{t-1},\qquad \sigma^2=2 .
\]
Compute the first three $\psi$-weights $\psi_0,\psi_1,\psi_2$ (from $\Theta(B)/\Phi(B)$), and give the point forecasts $X^{\,n}_{n+1}$ and $X^{\,n}_{n+2}$ in terms of $X_n$ and the residual $\hat\eps_n$. \hfill(1.5 pts)

\textbf{(c)} Report the forecast-error variances $P^{\,n}_{n+1},P^{\,n}_{n+2},P^{\,n}_{n+3}$ as exact numbers, and explain in one sentence the value $\acvs_0$ they approach as $h\to\infty$ (you may use $\acvs_0=\sigma^2(1+2\phi\theta+\theta^2)/(1-\phi^2)$ for an ARMA(1,1)). \hfill(1.5 pts)
\statementend

\textbf{(a)} Consider any linear predictor $X^{\,n}_{n+h}=\sum_{i=1}^n\beta_iX_i$. Substituting the linear representation $X_i=\sum_{j=0}^\infty\psi_j\eps_{i-j}$,
\[
X^{\,n}_{n+h}=\sum_{i=1}^n\beta_i\sum_{j=0}^\infty\psi_j\eps_{i-j}=\sum_{j=0}^\infty c_j\,\eps_{n-j}
\]
for some coefficients $c_j$: any predictor built from $X_1,\dots,X_n$ is a linear combination of the innovations up to time $n$ (it ``starts at $\eps_n$''). Now write $X_{n+h}=\sum_{j=0}^\infty\psi_j\eps_{n+h-j}$ and split off the future innovations $\eps_{n+1},\dots,\eps_{n+h}$ (indices $j=0,\dots,h-1$). Since the $\eps$'s are uncorrelated with common variance $\sigma^2$, the cross terms vanish and
\[
\begin{aligned}
\E\bigl[(X_{n+h}-X^{\,n}_{n+h})^2\bigr]
&=\E\Bigl[\Bigl(\sum_{j=0}^\infty\psi_j\eps_{n+h-j}-\sum_{j=0}^\infty c_j\eps_{n-j}\Bigr)^2\Bigr]\\
&=\sigma^2\Bigl(\sum_{j=0}^{h-1}\psi_j^2+\sum_{j=h}^\infty(\psi_j-c_{j-h})^2\Bigr).
\end{aligned}
\]
The first sum collects the future innovations, which no past-based predictor can touch (it is \emph{fixed}); the second sum matches $\psi_j$ against $c_{j-h}$ for the observable innovations. The second sum is minimised --- to $0$ --- by choosing $c_{j-h}=\psi_j$, i.e.\ $c_j=\psi_{j+h}$. The predictor is then
\[
X^{\,n}_{n+h}=\sum_{j=0}^\infty\psi_{j+h}\eps_{n-j}=\sum_{j=h}^\infty\psi_j\eps_{n+h-j},
\]
the residual is $X_{n+h}-X^{\,n}_{n+h}=\sum_{j=0}^{h-1}\psi_j\eps_{n+h-j}$, and the minimal MSE is
\[
P^{\,n}_{n+h}=\sigma^2\sum_{j=0}^{h-1}\psi_j^2 .
\]
(Operationally: \emph{write the future and erase the future innovations.})

\textbf{(b)} Here $\phi=\tfrac12$ and $\theta=\tfrac{3}{10}$ (with $\Phi(B)=1-\tfrac12B$, $\Theta(B)=1+\tfrac{3}{10}B$). The $\psi$-weights come from $\Theta(B)/\Phi(B)=\sum_j\psi_jB^j$:
\[
\frac{1+\tfrac{3}{10}B}{1-\tfrac12 B}
=\Bigl(1+\tfrac{3}{10}B\Bigr)\sum_{k=0}^\infty\Bigl(\tfrac12\Bigr)^k B^k .
\]
Matching powers of $B$,
\[
\psi_0=1,\qquad
\psi_1=\phi+\theta=\tfrac12+\tfrac{3}{10}=\tfrac{4}{5},\qquad
\psi_2=\phi\,\psi_1=\tfrac12\cdot\tfrac45=\tfrac25
\]
(in general $\psi_j=(\phi+\theta)\phi^{\,j-1}$ for $j\ge1$). The point forecasts are obtained by zeroing future innovations in $X_{n+h}=\tfrac12X_{n+h-1}+\eps_{n+h}+\tfrac{3}{10}\eps_{n+h-1}$:
\[
\begin{aligned}
X^{\,n}_{n+1}&=\tfrac12 X_n+\underbrace{\eps_{n+1}}_{\to0}+\tfrac{3}{10}\hat\eps_n
=\tfrac12 X_n+\tfrac{3}{10}\hat\eps_n,\\[2pt]
X^{\,n}_{n+2}&=\tfrac12 X^{\,n}_{n+1}+\underbrace{\eps_{n+2}}_{\to0}+\tfrac{3}{10}\underbrace{\eps_{n+1}}_{\to0}
=\tfrac12 X^{\,n}_{n+1}=\tfrac14 X_n+\tfrac{3}{20}\hat\eps_n .
\end{aligned}
\]
(For $h\ge2$ the MA term has vanished and the forecast obeys $X^{\,n}_{n+h}=\tfrac12 X^{\,n}_{n+h-1}$, decaying geometrically to the mean $0$.)

\textbf{(c)} With $\sigma^2=2$ and $\psi_0=1,\psi_1=\tfrac45,\psi_2=\tfrac25$, the forecast-error variances $P^{\,n}_{n+h}=\sigma^2\sum_{j=0}^{h-1}\psi_j^2$ are
\[
\begin{aligned}
P^{\,n}_{n+1}&=\sigma^2\,\psi_0^2=2\cdot1=2,\\
P^{\,n}_{n+2}&=\sigma^2\bigl(\psi_0^2+\psi_1^2\bigr)=2\Bigl(1+\tfrac{16}{25}\Bigr)=2\cdot\tfrac{41}{25}=\tfrac{82}{25}=3.28,\\
P^{\,n}_{n+3}&=\sigma^2\bigl(\psi_0^2+\psi_1^2+\psi_2^2\bigr)=2\Bigl(1+\tfrac{16}{25}+\tfrac{4}{25}\Bigr)=2\cdot\tfrac{45}{25}=\tfrac{18}{5}=3.6 .
\end{aligned}
\]
These increase toward the marginal variance
\[
\acvs_0=\sigma^2\,\frac{1+2\phi\theta+\theta^2}{1-\phi^2}
=2\cdot\frac{1+2\cdot\tfrac12\cdot\tfrac{3}{10}+\tfrac{9}{100}}{1-\tfrac14}
=2\cdot\frac{1+\tfrac{3}{10}+\tfrac{9}{100}}{\tfrac34}
=2\cdot\frac{139/100}{3/4}=\frac{278}{75}\approx3.707 ,
\]
which is exactly $\sigma^2\sum_{j=0}^\infty\psi_j^2$: as $h\to\infty$ the forecast loses all information and the error variance saturates at the unconditional variance $\acvs_0$ ($2<3.28<3.6<3.707$).

% ---------------------------------------------------------------
\solproblem{5}
\statementstart
\textbf{(Revision of Week 9: VAR stability and the diagonal VAR(1) covariance.)}

\textbf{(a)} Consider the bivariate $\mathrm{VAR}(2)$ process $X_t=\Phi_1X_{t-1}+\Phi_2X_{t-2}+\eps_t$ with
\[
\Phi_1=\begin{pmatrix}0.4 & 0\\[2pt] 0.5 & 0.9\end{pmatrix},\qquad
\Phi_2=\begin{pmatrix}0 & 0\\[2pt] 0.3 & -0.2\end{pmatrix}.
\]
Form the reverse characteristic polynomial $\det\bigl(I_2-\Phi_1 z-\Phi_2 z^2\bigr)$, exploit the lower-triangular structure to factor it, find all its roots, and decide whether the VAR(2) is stable. State the companion-matrix ($F$) form of the stability test and explain why it gives the same conclusion. \hfill(3 pts)

\textbf{(b)} Now take the \emph{diagonal} $\mathrm{VAR}(1)$ process $X_t=\Phi\,X_{t-1}+\eps_t$ with
\[
\Phi=\diag\!\bigl(\tfrac12,\tfrac14\bigr),\qquad
\{\eps_t\}\sim\mathrm{WN}(0,I_2).
\]
Write its $\mathrm{MA}(\infty)$ representation and compute the stationary covariance $\Gamma_0$ and the lag-$1$ matrix $\Gamma_1$ in closed form (exact fractions). \hfill(1.5 pts)

\textbf{(c)} For the process of part (b), what is the coherence $r_{1,2}(f)$ between the two channels at every frequency, and why? \hfill(0.5 pts)
\statementend

\textbf{(a)} Build $M(z)=I_2-\Phi_1 z-\Phi_2 z^2$:
\[
M(z)=\begin{pmatrix}
1-0.4z & 0\\[2pt]
-0.5z-0.3z^2 & 1-0.9z+0.2z^2
\end{pmatrix}.
\]
This matrix is lower triangular, so its determinant is the product of the diagonal entries (the off-diagonal terms $\Phi_1^{(2,1)},\Phi_2^{(2,1)}$ do \emph{not} enter):
\[
\det M(z)=(1-0.4z)\bigl(1-0.9z+0.2z^2\bigr).
\]
\emph{Roots.} The first factor gives $z=1/0.4=2.5$. The quadratic $0.2z^2-0.9z+1=0$ has
\[
z=\frac{0.9\pm\sqrt{0.81-0.8}}{0.4}=\frac{0.9\pm0.1}{0.4}
\ \Longrightarrow\ z=2.5\ \text{ or }\ z=2 .
\]
So the three roots are $z\in\{2.5,\,2.5,\,2\}$, all with $\abs z>1$ (strictly outside the unit circle). Hence the $\mathrm{VAR}(2)$ is \textbf{stable} (and therefore stationary). \emph{Sanity check:} $\det M(0)=\det I_2=1$, as required.

\emph{Companion form.} Stacking $Y_t=(X_t\Tr,X_{t-1}\Tr)\Tr\in\Reals^4$ turns the VAR(2) into the VAR(1) $Y_t=FY_{t-1}+U_t$ with
\[
F=\begin{pmatrix}\Phi_1 & \Phi_2\\ I_2 & 0\end{pmatrix}
=\begin{pmatrix}
0.4 & 0 & 0 & 0\\
0.5 & 0.9 & 0.3 & -0.2\\
1 & 0 & 0 & 0\\
0 & 1 & 0 & 0
\end{pmatrix}.
\]
The companion test is ``all eigenvalues of $F$ strictly inside the unit circle''. It is the reciprocal of the root test, because $\det(I_2-\Phi_1 z-\Phi_2 z^2)=z^{4}\det(z^{-1}I_4-F)$ up to a constant, so the eigenvalues of $F$ are the reciprocals $1/z_i$ of the roots. Here $\lambda=1/2.5=0.4$, $1/2.5=0.4$, $1/2=0.5$ (and a structural $\lambda=0$ from the rank-deficient $\Phi_2$), all with $\abs\lambda<1$: \emph{same} conclusion, stable.

\textbf{(b)} Since $\Phi=\diag(\tfrac12,\tfrac14)$ has spectral radius $\tfrac12<1$, the VAR(1) is stable and back-substitution gives the matrix geometric ($\mathrm{MA}(\infty)$) filter
\[
X_t=\sum_{j=0}^{\infty}\Phi^{\,j}\eps_{t-j},
\qquad
\Phi^{\,j}=\diag\!\bigl(2^{-j},\,4^{-j}\bigr).
\]
By the VAR(1) covariance formula $\Gamma_0=\sum_{k\ge0}\Phi^k\Sigma(\Phi^k)\Tr$ with $\Sigma=I_2$ (everything diagonal, so the channels decouple into independent AR(1)'s),
\[
\Gamma_0=\sum_{k=0}^{\infty}\diag\!\bigl(4^{-k},\,16^{-k}\bigr)
=\diag\!\left(\frac{1}{1-\tfrac14},\ \frac{1}{1-\tfrac{1}{16}}\right)
=\boxed{\ \diag\!\left(\tfrac43,\ \tfrac{16}{15}\right)\ }.
\]
(Equivalently, the Lyapunov equation $\Gamma_0=\Phi\Gamma_0\Phi\Tr+I_2$ gives componentwise $g_{jj}=\phi_j^2 g_{jj}+1\Rightarrow g_{jj}=1/(1-\phi_j^2)$, i.e.\ $1/(1-\tfrac14)=\tfrac43$ and $1/(1-\tfrac{1}{16})=\tfrac{16}{15}$.) The lag-$1$ matrix follows from Yule--Walker $\Gamma_1=\Phi\Gamma_0$:
\[
\Gamma_1=\diag\!\bigl(\tfrac12,\tfrac14\bigr)\diag\!\bigl(\tfrac43,\tfrac{16}{15}\bigr)
=\boxed{\ \diag\!\left(\tfrac23,\ \tfrac{4}{15}\right)\ }.
\]

\textbf{(c)} Both $\Phi$ and $\Sigma=I_2$ are diagonal, so by the formula in (b) \emph{every} $\Gamma_\tau$ is diagonal: all off-diagonal cross-covariances vanish at every lag. Hence the cross-spectrum $S_{1,2}(f)=\sum_\tau\acvs_\tau^{(1,2)}e^{-2\pi if\tau}\equiv0$, and the coherence
\[
r_{1,2}(f)=\frac{S_{1,2}(f)}{\sqrt{S_{1,1}(f)S_{2,2}(f)}}=0\quad\text{for all }f .
\]
The two channels are completely uncoupled: the off-diagonal entries of $\Phi$ (here zero) are exactly the cross-channel dynamics, and there is no shared innovation ($\Sigma$ diagonal), so nothing links the series at any frequency.

\end{document}
